\documentclass[10pt,a4paper]{amsart}
\usepackage[margin=19mm]{geometry}
\usepackage{amsmath,amssymb,mathtools}
\usepackage[scr=rsfs,cal=euler]{mathalfa}
\usepackage{xfrac}
\usepackage[unicode,hidelinks,bookmarks=false]{hyperref}
\newcommand{\ie}{i.e.\ }
\newcommand{\cf}{cf.\ }
\newcommand{\resp}{respectively\ }
\newcommand{\Subsection}{\subsection}
\providecommand{\nobreakdash}{\nobreak\hbox{-}}
\setlength{\emergencystretch}{2em}
\multlinegap0pt
\usepackage{xcolor}
\usepackage{amssymb}
\usepackage{mathtools}
\usepackage{nicefrac}
\newtheorem{lemma}{Lemma}
\newcommand{\PP}{\mathbb{P}}
\title{Sharp large deviation estimates for heavy-tailed extrema}
\author{Jos\'e M. Zapata}
\date{Source: arXiv:2512.24352v2 (CC BY 4.0)}
\begin{document}
\maketitle

\noindent\emph{Source and licence.} Sharp large deviation estimates for heavy-tailed extrema. Version arXiv:2512.24352v2 (CC BY 4.0), distributed by arXiv under the Creative Commons Attribution 4.0 International licence, http://creativecommons.org/licenses/by/4.0/, as recorded in the arXiv licence metadata for this exact version and shown on the abstract page https://arxiv.org/abs/2512.24352v2. Full licence notice: \texttt{licenses/arxiv-2512.24352-CC-BY-4.0.txt}.\par\smallskip

\noindent\textit{Editorial scope.} Counted unit: the article's lemma lem:limsup (source0.tex lines 629--635). The article's complete original proof of it is delivered with it (source0.tex lines 637--678, one \texttt{proof} environment of 42 lines). No mathematics is written, repaired or re-derived: the statement and the proof are the article's own lines, in the article's own notation, and no retained line is substituted or re-flowed.  Retained uncounted as the source-local context the proof needs: lines 613--623.  Nothing was omitted from any retained span, so the package carries no omission record.  No retained line contains a citation, so the wrapper carries no bibliography and no bibliography entry.  No retained line contains a figure environment, an image-inclusion command or drawing code, so no image is delivered and the package declares no asset dependency.  Editorial formatting only, in addition: an AMS \texttt{amsart} preamble that restates the article's own declarations and macros used by the retained lines, namely the environments \texttt{lemma} on separate counters, together with the standard packages those lines need.  The retained environments print in this document as Lemma 1 in order of appearance, whereas the article prints the same items with the numbering of its own full document.  The complete native source of the article as distributed in the arXiv e-print for this exact version is bundled unchanged as \texttt{sources/191885/source0.tex}.
\begin{lemma}\label{lem:Potter} 
For every $\varepsilon>0$ there exists $t_0>0$ such that
\[
(1-\varepsilon)\,x^{-\alpha-\varepsilon}
\;\le\;
\frac{\bar F(tx)}{\bar F(t)}
\;\le\;
(1+\varepsilon)\,x^{-\alpha+\varepsilon},
\]
for all $t\ge t_0$ and all $x\ge x_0$.
\end{lemma}

\begin{lemma}\label{lem:limsup}
For every \(x \ge 1\),
\[
\limsup_{n\to\infty}\frac{1}{\log n}\log \PP(Z_n>x)\le -I(x).
\]
Recall that $I(x)=\log x$.
\end{lemma}

\begin{proof}
Fix \(x\ge1\).
Using the elementary identity
\[
1-u^n=(1+u+\cdots+u^{n-1})(1-u), \qquad u\in[0,1],
\]
we obtain the inequality
\[
1-u^n \le n(1-u), \qquad u\in[0,1].
\]
Applying this with \(u=F(t_n(x))\), we obtain
\[
\PP(Z_n>x)
=
1-F(t_n(x))^n
\le
n\,\bar F(t_n(x)).
\]

Since \(x\ge1\), we have \(t_n(x)\ge a_n\to\infty\) as \(n\to\infty\).
Fix \(\varepsilon>0\). By the Potter bounds (Lemma~\ref{lem:Potter}), for all \(n\) sufficiently large,
\[
\bar F(t_n(x))
\le
(1+\varepsilon)\,\bar F(a_n)\,x^{\frac{\log n}{\alpha}(-\alpha+\varepsilon)}.
\]
Since \(\bar F(a_n)\le 1/n\), we deduce
\[
\PP(Z_n>x)
\le
(1+\varepsilon)\,x^{\frac{\log n}{\alpha}(-\alpha+\varepsilon)}.
\]

Taking logarithms, dividing by \(\log n\) , and taking the limit superior yield 
\[
\limsup_{n\to\infty}
\frac{1}{\log n}\log \PP(Z_n>x)
\le
\left(-1+\frac{\varepsilon}{\alpha}\right)\log x.
\]
Letting \(\varepsilon\downarrow0\) completes the proof.
\end{proof}

\end{document}
